\section{Six Hypotheses Refuted, Three Claims Withdrawn}
\label{sec:negative}

We set out to find regimes where the diffusion model beats the autoregressive
baseline, and generated hypotheses accordingly. Six were refuted and three
earlier claims were withdrawn. They are reported here at the same length and
in the same register as the confirmations, because a paper whose central
methodological claim is ``price the baseline before you compare against it''
has exactly one asset, and that asset is spent the moment a negative result
gets softened into a footnote.

\subsection{Six refuted hypotheses}

\begin{table}[h]
\centering
\caption{Six hypotheses of ours that did not survive measurement.}
\label{tab:refuted}
\begin{tabularx}{\linewidth}{@{}p{4.6cm}L@{}}
\toprule
\thead{Hypothesis} & \thead{What we measured instead} \\
\midrule
\tabfont Joint denoising resists input perturbation
  & \tabfont The diffusion model degraded more in \textbf{10 of 10} categories, by up to $8.2\times$ the autoregressive drop. Discussed below. \\
\tabfont Bidirectional conditioning wins at infilling
  & \tabfont \ARthree\ scores $0.993$--$0.999$ across all suffix sizes; the diffusion model holds to $0.960$ then falls to $0.753$ at $S{=}8$. Discussed below. \\
\tabfont Diffusion has a refinement mechanism that pays off
  & \tabfont It has the mechanism --- selective re-masking gains $+0.023$ F1 over three rounds --- but after three rounds it still trails the autoregressive model after one. The mechanism is real; the win is not. \\
\tabfont Diffusion beats autoregression under bounded deadlines
  & \tabfont This was defect 4. The autoregressive model does complete within the deadline; the metric had been counting valid falsy values as missing keys (\S\ref{sec:setup}). \\
\tabfont Wide-schema cost scales as $O(F^2)$ for diffusion
  & \tabfont Never measured. The consistency metric is three hardcoded rules over four fields at \emph{every} schema width, so the curve it produces cannot show quadratic structure in $F$. The framing is dropped. \\
\tabfont Quality is invariant to output length
  & \tabfont Quality rises steeply with $L$ and then plateaus (\S\ref{sec:exchange}). The apparent invariance came from the contaminated 12-field metric. \\
\bottomrule
\end{tabularx}
\end{table}

\paragraph{Perturbation robustness, refuted in every category}
We predicted that denoising a whole canvas jointly would make the diffusion
model more tolerant of noisy input than left-to-right decoding. Across ten
perturbation categories at full intensity, it was worse in all ten. Two
categories --- whitespace and abbreviation --- left the autoregressive model
completely flat while costing the diffusion model $0.081$ and $0.005$ F1. The
largest ratio is separator changes, at $8.2\times$. One hypothesis for the
direction: left-to-right conditioning propagates input noise forward only,
while bidirectional conditioning also propagates it backward. We did not test
that, and it stays a hypothesis.

\paragraph{Infilling, refuted for a reason worth stating}
The prediction was structural rather than statistical: with a fixed right
context, suffix tokens are never in the denoised set, so the diffusion model
should preserve them byte-identically by construction, while an autoregressive
model must regenerate them. The experiment ($S \in \{0,2,4,8\}$ suffix fields,
$n{=}50$) refuted it, but the mechanism matters --- the autoregressive model
scores $0.993$--$0.999$ because it trivially regenerates the whole record from
the prompt. This task is copy-heavy, and copying is something autoregressive
decoding does well. The structural argument may still hold where the suffix is
long and cannot be regenerated from context, such as code infilling. It does
not hold on short JSON fields, which is what we measured, so it is not a claim
we make.

\subsection{Three withdrawn claims}

\begin{table}[h]
\centering
\caption{Three claims that appeared in earlier drafts and are cut from this
one. Withdrawn claims are removed.}
\label{tab:withdrawn}
\begin{tabularx}{\linewidth}{@{}p{4.6cm}L@{}}
\toprule
\thead{Claim} & \thead{Why it is gone} \\
\midrule
\tabfont A classical parser matches LLM extraction
  & \tabfont Field-F1 $0.68$ against $0.85$; the gap is too large. Worse, our ``Drain''~\cite{drain} was a regex fallback rather than the reference implementation, and the fallback and the gold generator are the same code, so on semantic fields it was grading itself. \\
\tabfont Structural and semantic fields separate cleanly
  & \tabfont The semantic gold labels are placeholder output of that same parser, and one of them is a constant. ``Semantic F1'' may be measuring agreement with a regex. \\
\tabfont An M/D/1 queue predicts diffusion latency under load
  & \tabfont It matches the measured saturation knee under Poisson arrivals and stops being predictive under bursty (MMPP) arrivals, which is what log streams actually look like. \\
\bottomrule
\end{tabularx}
\end{table}

\paragraph{A pattern in the refutations}
Every refuted hypothesis had, at some earlier point, appeared to be confirmed.
In each case the apparent confirmation traced to an artifact of our own
harness rather than to a property of either architecture: the bounded-deadline
result was the completeness metric, the length-invariance result was the
contaminated metric, and the perturbation and refinement arms ran a
mis-provenanced model. The measurement discipline of \S\ref{sec:setup} is not
a preamble to the results. It is what separates these six from the ones we
kept.
